% =============================================================================
%  Sentinel DV — Technical White Paper & Implementation Notes
%  DVCon USA / DAC 2026 Companion Document
% =============================================================================
\documentclass[11pt,a4paper]{article}

\usepackage[T1]{fontenc}
\usepackage[utf8]{inputenc}
\usepackage{amsmath,amssymb}
\usepackage{graphicx}
\usepackage{xcolor}
\usepackage{booktabs}
\usepackage{listings}
\usepackage{microtype}
\usepackage{hyperref}
\usepackage[margin=1.0in]{geometry}
\usepackage{fancyhdr}
\usepackage{titlesec}
\usepackage{enumitem}
\usepackage{cite}
\usepackage{url}
\usepackage{caption}
\usepackage{parskip}

% Colors
\definecolor{codeblue}{RGB}{26,86,219}
\definecolor{codegray}{RGB}{107,114,128}
\definecolor{codegreen}{RGB}{5,122,85}
\definecolor{codered}{RGB}{180,0,0}
\definecolor{codeteal}{RGB}{6,148,162}
\definecolor{codedark}{RGB}{17,24,39}

\lstdefinestyle{yaml}{
  backgroundcolor=\color{codedark!95!white},
  basicstyle=\ttfamily\small\color{white},
  keywordstyle=\color{codeteal}\bfseries,
  stringstyle=\color{codegreen},
  commentstyle=\color{codegray}\itshape,
  breaklines=true, frame=none, numbers=left,
  numberstyle=\tiny\color{codegray},
  xleftmargin=1.5em,
}

\tcbset{codebox/.style={
  colback=codedark!95!white, colframe=codeblue,
  colupper=white, fontupper=\ttfamily\small,
  boxrule=0.5pt, arc=3pt,
}}

\hypersetup{colorlinks=true, linkcolor=codeblue, citecolor=codegreen, urlcolor=codeblue}

\titleformat{\section}{\large\bfseries\color{codeblue}}{{\arabic{section}.}}{0.5em}{}[\titlerule]
\titleformat{\subsection}{\normalsize\bfseries\color{codedark}}{{\arabic{section}.\arabic{subsection}}}{0.4em}{}
\titleformat{\subsubsection}{\normalsize\itshape}{}{0em}{}

\pagestyle{fancy}
\fancyhf{}
\fancyhead[L]{\small\textit{Sentinel DV Technical Notes}}
\fancyhead[R]{\small\textit{DVCon/DAC 2026}}
\fancyfoot[C]{\thepage}

% =============================================================================
\title{%
  {\color{codeblue}\textbf{Sentinel DV}}\\[6pt]
  \large Technical White Paper \& Implementation Notes\\[4pt]
  \normalsize DVCon USA / DAC 2026 — Companion Document
}
\author{Open-Source Contribution \\ \texttt{\href{https://github.com/kiranreddi/sentinel-dv}{github.com/kiranreddi/sentinel-dv}}}
\date{2026 \\ MIT License \quad \texttt{pip install sentinel-dv}}

% =============================================================================
\begin{document}
\maketitle
\thispagestyle{fancy}

\begin{abstract}
This document provides detailed technical notes to accompany the Sentinel DV conference
paper~\cite{sdv_paper}. It covers installation, configuration, tool API reference,
adapter internals, intelligence algorithm details, and advanced deployment patterns.
Intended for DV engineers evaluating or extending the framework.
\end{abstract}

\tableofcontents
\newpage

% =============================================================================
\section{Installation \& Quick Start}

\subsection{Requirements}
\begin{itemize}
  \item Python 3.10, 3.11, or 3.12
  \item No EDA tool licenses required (parsers consume log/HTML artifacts only)
  \item Disk space: $\sim$50\,MB base, plus artifact storage
  \item Network: none (all local; optional PyPI for install)
\end{itemize}

\subsection{Install}
\begin{lstlisting}[style=yaml,numbers=none,xleftmargin=0em]
pip install sentinel-dv
\end{lstlisting}

Or from source:
\begin{lstlisting}[style=yaml,numbers=none,xleftmargin=0em]
git clone https://github.com/kiranreddi/sentinel-dv\\
cd sentinel-dv\\
pip install -e ".[dev]"
\end{lstlisting}

\subsection{Minimal Configuration}

Create \texttt{sentinel\_dv.yaml} in your project root:

\begin{lstlisting}[style=yaml,caption={Minimal sentinel\_dv.yaml}]
# Sentinel DV configuration
artifact_roots:
  - /path/to/sim/logs          # directory containing *.log UVM logs
  - /path/to/coverage/report   # directory containing dashboard.html
  - /path/to/regression/xmls   # directory containing results.xml

db_path: ./sentinel_dv.db
\end{lstlisting}

\subsection{Start the MCP Server}
\begin{lstlisting}[style=yaml,numbers=none,xleftmargin=0em]
sentinel-dv serve --config sentinel\_dv.yaml
\end{lstlisting}

The server binds to \texttt{stdio} transport by default (required for Claude/VS~Code
Copilot integration). Add \texttt{--transport sse --port 8765} for network access.

\subsection{Connect an AI Client}

For Claude Desktop (\texttt{\textasciitilde/.config/claude/claude\_desktop\_config.json}):
\begin{lstlisting}[style=yaml,caption={Claude Desktop MCP config}]
{
  "mcpServers": {
    "sentinel-dv": {
      "command": "sentinel-dv",
      "args": ["serve", "--config", "sentinel_dv.yaml"]
    }
  }
}
\end{lstlisting}

For VS Code Copilot (\texttt{.vscode/mcp.json} in workspace):
\begin{lstlisting}[style=yaml,caption={VS Code Copilot MCP config}]
{
  "inputs": [],
  "servers": {
    "sentinel-dv": {
      "type": "stdio",
      "command": "sentinel-dv",
      "args": ["serve", "--config", "sentinel_dv.yaml"]
    }
  }
}
\end{lstlisting}

% =============================================================================
\section{Configuration Reference}

\begin{lstlisting}[style=yaml,caption={Full sentinel\_dv.yaml reference}]
# === Artifact Discovery ===
artifact_roots:
  - /path/to/primary/artifacts
  - /path/to/coverage/reports

# === Database ===
db_path: ./sentinel_dv.db       # SQLite database path
auto_index: true                 # Re-index on server start

# === Coverage ===
coverage:
  html_glob_patterns:           # Glob patterns to find coverage HTML
    - "*/coverage/report/dashboard.html"    # VCS URG
    - "*/html_report/overalldu.js"          # Questa vcover
    - "*/coverage/index.html"              # Generic
  vcs_score_threshold: 75.0    # Warn if total < threshold
  questa_score_threshold: 75.0

# === Regression Health ===
health:
  pass_weight:      0.30
  coverage_weight:  0.35
  assertion_weight: 0.15
  flakiness_weight: 0.10
  cross_sim_weight: 0.10
  target_coverage:  75.0        # %  used in health computation

# === Path Redaction ===
redact_paths: true
redact_patterns:                # Additional regex patterns to redact
  - '/home/[^/]+'
  - '/workarea/[^/]+'

# === Simulator Support ===
simulators:
  - vcs
  - questa
  - xcelium
\end{lstlisting}

% =============================================================================
\section{MCP Tool API Reference}

All tools communicate via JSON-RPC 2.0. Inputs and outputs are Pydantic-validated.
Below is a condensed reference; full schemas are in
\texttt{sentinel\_dv/tools/schemas.py}.

\subsection{Run Tracking Tools}

\subsubsection{\texttt{runs.list}}
\textbf{Parameters:} \texttt{page} (int, default 1), \texttt{page\_size} (int, default 20),
\texttt{simulator} (str, optional), \texttt{status} (str: pass/fail/all).\\
\textbf{Returns:} \texttt{runs} (list), \texttt{total} (int), \texttt{page\_info} (object).

\subsubsection{\texttt{runs.get}}
\textbf{Parameters:} \texttt{run\_id} (str, required).\\
\textbf{Returns:} Full run object under \texttt{run} key including test counts, timestamps,
simulator, duration, log excerpts.

\subsubsection{\texttt{runs.diff}}
\textbf{Parameters:} \texttt{run\_id\_a}, \texttt{run\_id\_b} (str).\\
\textbf{Returns:} \texttt{added\_failures}, \texttt{fixed\_tests}, \texttt{coverage\_delta}.

\subsubsection{\texttt{runs.cross\_sim}}
Compares the same test suite across VCS, Questa, and Xcelium.\\
\textbf{Returns:} \texttt{simulators\_compared}, \texttt{divergent\_tests} list,
\texttt{consistency\_score} (0--100).

\subsection{Coverage Intelligence Tools}

\subsubsection{\texttt{coverage.summary}}
\textbf{Parameters:} \texttt{run\_id} (str, optional — uses latest if omitted).\\
\textbf{Returns:} \texttt{summaries} list, each with \texttt{kind} (code/functional/assertion/plan),
\texttt{metrics} list of \{name, covered, total\}.

\textbf{Important:} Returns \texttt{summaries} (list), not \texttt{summary} (singular).

\subsubsection{\texttt{coverage.trend}}
\textbf{Parameters:} \texttt{window} (int, default 10 runs), \texttt{metric} (str, default "total").\\
\textbf{Returns:} \texttt{data\_points} list, \texttt{summary.direction} (improving/declining/stable).

\textbf{Important:} Direction is under \texttt{summary.direction}, not top-level.

\subsubsection{\texttt{coverage.gaps}}
\textbf{Parameters:} \texttt{threshold} (float, default 25.0 \%).\\
\textbf{Returns:} \texttt{gaps} list with \texttt{scope}, \texttt{kind}, \texttt{covered\_pct}.

\subsubsection{\texttt{coverage.advisor}}
\textbf{Parameters:} \texttt{gap\_threshold} (float, default 25.0).\\
\textbf{Returns:} \texttt{advisories} list with \texttt{priority}, \texttt{sv\_constraint},
\texttt{uvm\_sequence}, \texttt{recommendation}.

\subsection{Regression Health Tool}

\subsubsection{\texttt{regression.health}}
No required parameters. Returns:
\begin{lstlisting}[style=yaml]
{
  "health_score": 77.5,          // 0-100 composite
  "band": "minor-issues",        // sign-off-ready | minor-issues |
                                 // coverage-gaps  | not-ready
  "grade": "B+",
  "components": {
    "pass_rate":    94.2,
    "coverage":     51.3,
    "assertions":   100.0,
    "flakiness":    98.1,
    "cross_sim":    87.5
  },
  "recommendations": [...]       // ranked list of action items
}
\end{lstlisting}

% =============================================================================
\section{Adapter Layer Internals}

\subsection{Multi-Simulator UVM Log Parser}

File: \texttt{sentinel\_dv/adapters/uvm\_log.py}

The parser uses five regex patterns, tried in priority order per log line:

\begin{enumerate}
  \item \textbf{VCS native}: \verb|UVM_ERROR file.sv(88) @ 500ns: comp [CAT] msg|
  \item \textbf{Questa}: \verb|# UVM_ERROR @ 500 ns: comp.path [CAT] msg|
  \item \textbf{Xcelium}: \verb|xcelium: UVM_ERROR @ 500ns: comp [CAT] msg|
  \item \textbf{Generic UVM}: matches any remaining \verb|UVM_(ERROR|FATAL)| line
  \item \textbf{Count-summary skip}: \verb|UVM_(ERROR|FATAL)\s*:\s*\d+| — \textbf{skipped}
\end{enumerate}

\textbf{False-positive elimination}: Pattern 5 matches the summary lines that VCS
appends to every log (e.g., \texttt{UVM\_ERROR :\ \ \ 0}). Without this skip rule,
count-only report summaries are easily misclassified as real failures.

\subsection{VCS URG HTML Coverage Parser}

File: \texttt{sentinel\_dv/adapters/coverage\_reports.py} — \texttt{\_parse\_urg\_html()}

Algorithm:
\begin{enumerate}
  \item Read \texttt{dashboard.html}; content-sniff for ``Total Coverage Summary''
  \item Locate the summary table via regex \verb|<table[^>]*>.*?</table>|
  \item For each \verb|<tr>|, extract metric name and percentage value
  \item Map vendor column names: SCORE→total, LINE→line, TOGGLE→toggle, etc.
  \item Return \texttt{CoverageSummary} Pydantic object
\end{enumerate}

\textbf{Handled edge cases}: numeric locale formatting (\texttt{87.30} vs
\texttt{87,30}), N/A values, missing columns (old URG versions omit verification
plan).

\subsection{Questa vcover JS Coverage Parser}

File: \texttt{sentinel\_dv/adapters/coverage\_reports.py} — \texttt{\_parse\_questa\_overalldu()}

The \texttt{overalldu.js} file is generated as a single-line JavaScript assignment:
\begin{lstlisting}[style=yaml,numbers=none,xleftmargin=0em]
var overallDu = \{``items'': [\{``type'': ``stmt'', ``covered'': 1024, ``total'': 1200, ...\}]\};
\end{lstlisting}

Standard regex fails on deeply nested structures. The parser uses a
\textbf{balanced-brace scanner}:
\begin{enumerate}
  \item Locate the opening \texttt{\{} after the assignment
  \item Walk characters, tracking brace depth; collect substring until depth returns to 0
  \item Parse the extracted substring as JSON (after minor JS→JSON normalization)
\end{enumerate}

\textbf{Extracted fields}: stmt, branch, toggle, fsm, assertion, functional, total.
Each yields \texttt{hit} and \texttt{total} counts for percentage computation.

% =============================================================================
\section{DV Intelligence Algorithms}

\subsection{Regression Health Score}

The composite health score $H \in [0, 100]$ is:
\[
H = w_P \cdot P + w_C \cdot C + w_A \cdot A + w_F \cdot (1 - F) + w_X \cdot X
\]

where weights $(w_P, w_C, w_A, w_F, w_X) = (0.30, 0.35, 0.15, 0.10, 0.10)$ are
configurable via \texttt{sentinel\_dv.yaml}. Components:
\begin{itemize}
  \item $P$: pass rate = passing tests / total tests (0--100)
  \item $C$: coverage score from \texttt{coverage.summary} total metric (0--100)
  \item $A$: assertion health from indexed assertion definitions/status. If no
        assertion evidence is indexed, this component is marked unavailable and
        the remaining weights are renormalized.
  \item $F$: flakiness fraction = tests that diverge across simulator runs (0--1)
  \item $X$: cross-simulator consistency from \texttt{runs.cross\_sim} (0--100)
\end{itemize}

Band thresholds are configurable. Default: 90/75/55 for sign-off/minor/gap bands.

\subsection{Coverage Trend}

For a window of $N$ consecutive coverage reports ordered by timestamp, the trend
direction is determined by a simple linear regression slope on the total coverage
percentage series. Threshold: $|\text{slope}| < 0.5\%$/run is classified as
\texttt{stable}; positive slope as \texttt{improving}; negative as \texttt{declining}.

\subsection{Test Clustering}

Failing test error messages are classified by string-matching heuristics:
\begin{itemize}
  \item \texttt{timeout}: contains ``timeout'', ``TIMEOUT'', ``time limit''
  \item \texttt{assertion}: UVM\_FATAL or SVA \texttt{ASSERT\_*} message
  \item \texttt{data\_mismatch}: contains ``MISMATCH'', ``Expected'', ``got''
  \item \texttt{uvm\_phase}: contains ``phase'', ``objection'', ``raise/drop''
  \item \texttt{unknown}: fallback category
\end{itemize}

Within each cluster, the representative failure is the most frequently occurring
message pattern. Cluster size ranks recommendations.

\subsection{Coverage Advisor Protocol Detection}

The advisor detects the underlying protocol from signal names in coverage gaps:
\begin{center}
\begin{tabular}{ll}
\toprule
\textbf{Pattern in scope name} & \textbf{Protocol} \\
\midrule
\verb|awaddr|, \verb|araddr|, \verb|awburst| & AXI4 \\
\verb|haddr|, \verb|htrans|, \verb|hburst| & AHB \\
\verb|paddr|, \verb|penable|, \verb|psel| & APB \\
\verb|txreq|, \verb|rxdat|, \verb|txrsp| & CHI \\
\verb|tl_a_|, \verb|tl_d_| & TileLink \\
\verb|pcie_|, \verb|gen3_| & PCIe \\
\bottomrule
\end{tabular}
\end{center}

% =============================================================================
\section{AXI4 UVM Demo}

The repository includes a complete AXI4-Lite verification demo in
\texttt{demo/axi4\_uvm/}.

\subsection{Demo Structure}
\begin{lstlisting}[style=yaml,numbers=none,xleftmargin=0em]
demo/axi4\_uvm/\\
  config.yaml           \# Sentinel DV config for the demo\\
  vcs/                  \# VCS simulator artifacts\\
    results.xml         \# JUnit-format test results\\
    assertions/         \# Assertion JSON files\\
    coverage/           \# dashboard.html (VCS URG)\\
    waveforms/          \# .wave.json waveform summaries\\
  questa/               \# Questa simulator artifacts (same layout)\\
  xcelium/              \# Xcelium simulator artifacts (same layout)\\
  README.md             \# How to run and use
\end{lstlisting}

\subsection{Test Cases}
Four test cases cover the AXI4-Lite verification plan:
\begin{enumerate}
  \item \texttt{axi4\_bk2bk\_test}: Back-to-back write/read transactions
  \item \texttt{axi4\_err\_resp\_test}: Error response checking (SLVERR/DECERR)
  \item \texttt{axi4\_boundary\_test}: 4KB boundary alignment
  \item \texttt{axi4\_stress\_test}: Random burst stress (reveals cross-sim divergence)
\end{enumerate}

\subsection{Running the Demo}
\begin{lstlisting}[style=yaml,numbers=none,xleftmargin=0em]
cd sentinel-dv\\
sentinel-dv index --config demo/axi4\_uvm/config.yaml\\
\# Then open Claude Desktop and ask:\\
\# "What's the health of the AXI4 demo regression?"
\end{lstlisting}

Expected output: health score $\sim$72, band \texttt{coverage-gaps},
3 divergent tests between Questa and VCS/Xcelium for \texttt{axi4\_stress\_test}.

% =============================================================================
\section{Extending Sentinel DV}

\subsection{Adding a New Adapter}

\begin{enumerate}
  \item Create \texttt{sentinel\_dv/adapters/my\_format.py}
  \item Implement \texttt{parse(path: str) -> List[Artifact]} function
  \item Register in \texttt{sentinel\_dv/adapters/\_\_init\_\_.py}
  \item Add glob patterns to \texttt{ARTIFACT\_GLOBS} in \texttt{indexer.py}
  \item Write unit tests in \texttt{tests/unit/test\_my\_format.py}
\end{enumerate}

\subsection{Adding a New MCP Tool}

\begin{enumerate}
  \item Add tool function to appropriate module in \texttt{sentinel\_dv/tools/}
  \item Decorate with \texttt{@mcp.tool()} from \texttt{fastmcp}
  \item Define Pydantic input/output schemas in \texttt{tools/schemas.py}
  \item Register in \texttt{sentinel\_dv/server.py}
  \item Add integration test in \texttt{tests/integration/}
\end{enumerate}

% =============================================================================
\section{CI/CD Integration}

\subsection{GitHub Actions}
\begin{lstlisting}[style=yaml,caption={.github/workflows/sentinel-dv.yml}]
name: Sentinel DV Analysis
on: [push]

jobs:
  dv-analysis:
    runs-on: ubuntu-latest
    steps:
      - uses: actions/checkout@v4
      - name: Install Sentinel DV
        run: pip install sentinel-dv
      - name: Download sim artifacts
        run: |
          # Copy/mount your simulation artifacts
          cp -r $ARTIFACT_DIR ./sim_artifacts
      - name: Run Sentinel DV health check
        run: |
          sentinel-dv index --config .sentinel_dv.yaml
          sentinel-dv check-health --min-score 70
\end{lstlisting}

\subsection{Jenkins Pipeline}
\begin{lstlisting}[style=yaml,caption={Jenkinsfile snippet}]
stage('DV Quality Gate') {
    steps {
        sh 'pip install sentinel-dv'
        sh 'sentinel-dv index --config sentinel_dv.yaml'
        sh 'sentinel-dv check-health --min-score ${params.MIN_HEALTH:-70}'
    }
}
\end{lstlisting}

% =============================================================================
\section{Troubleshooting}

\subsection{Common Issues}

\noindent\textbf{Coverage not found}: Ensure the HTML files exist at the paths in
\texttt{artifact\_roots}. Run \texttt{sentinel-dv list-artifacts} to see what
was discovered.

\noindent\textbf{0 runs indexed}: Check that \texttt{artifact\_roots} contains
directories with \texttt{*.log} or \texttt{results.xml} files. VCS logs must
contain \texttt{UVM\_ERROR} or \texttt{Errors:} lines to be detected as
simulation logs.

\noindent\textbf{Health score returns 0}: Run \texttt{regression.health} with
\texttt{--debug} flag or check that at least one run is indexed with
\texttt{runs.list}.

\noindent\textbf{coverage.summary returns empty summaries}: The HTML parser
requires the vendor-specific report format. Ensure you are using
\texttt{urg -report} (VCS) or \texttt{vcover report -html} (Questa).

% =============================================================================
\section{Future Roadmap}

\begin{itemize}
  \item \textbf{v2.2.0}: Direct binary database parsing (UCDB via \texttt{questa-vpm}, VDB via \texttt{vcs-dv})
  \item \textbf{v2.3.0}: FSDB/MXDB waveform signals and transitions
  \item \textbf{v3.0.0}: Formal property synthesis from coverage gap analysis
  \item \textbf{v3.1.0}: Cloud MCP endpoint + Docker image for CI/CD
  \item \textbf{v3.2.0}: Test generation from failure clusters using LLM fine-tuning
  \item \textbf{v4.0.0}: Testplan tracking with coverage closure metrics
\end{itemize}

\section*{References}

\begin{thebibliography}{9}
\bibitem{sdv_paper}
Sentinel DV Conference Paper, DVCon USA 2026.
\texttt{conference/paper/sentinel\_dv\_dvcon2026.pdf}

\bibitem{mcp2024}
Anthropic, ``Model Context Protocol,''
\url{https://modelcontextprotocol.io}, 2024.

\bibitem{fastmcp}
J.~Lowin, ``FastMCP,'' \url{https://github.com/jlowin/fastmcp}, 2024.
\end{thebibliography}

\end{document}
